\subsection{Research Directions and Evaluation}
% Source page: Deepfield.md Page 54; exactly one frame.
\setsource{}
\begin{frame}[c,label=md-54]{Six Research Directions}
\small
\renewcommand{\arraystretch}{1.15}
\begin{tabular*}{\textwidth}{@{\extracolsep{\fill}}L{0.35\textwidth}L{0.58\textwidth}@{}}
\toprule
\textbf{Direction} & \textbf{Research focus} \\
\midrule
\textbf{I. Knowledge and models} & Evidence, physical models, experience, and long-term state. \\[1.0ex]
\textbf{II. Hypotheses} & Hierarchical questions, mechanisms, cross-domain reasoning, evaluation. \\[1.0ex]
\textbf{III. Code and analysis} & Reproduction, reuse, development, analysis, simulation. \\[1.0ex]
\textbf{IV. Experiment design} & Feasible plans, settings, and observing resources. \\[1.0ex]
\textbf{V. Experimentation} & Beam, astrophysical, and numerical experiments. \\[1.0ex]
\textbf{VI. Benchmarking} & Literature, hypothesis, code, and experiment evaluation. \\[1.0ex]
\bottomrule
\end{tabular*}\par

\slidenote{Near term: a complete JUNA Loop. Longer term: transfer
across reactions, environments, and scientific fields.}
\end{frame}

% Source page: Deepfield.md Page 55; exactly one frame.
\setsource{}
\begin{frame}[t,label=md-55]{I. Knowledge-Guided Exploration and Model Discovery}
\small
\begin{itemize}
\setlength{\itemsep}{0.6em}
    \item \textbf{Evidence:} Papers, figures, formulas, code, EXFOR, ENSDF, and REACLIB / STARLIB, with provenance and validity.
    \item \textbf{Questions:} Candidate reactions and energy regions, gaps, conflicts, and inventories.
    \item \textbf{Models:} Compare mechanisms, diagnose mismatch, and explore constrained equation or program discovery.
    \item \textbf{Experience and state:} Retain configurations, anomalies, failures, feedback, hypothesis branches, and checkpoints.
\end{itemize}

\medskip
\textbf{Test:} Does continual evidence alignment improve quantitative decisions?
\smallskip

\textbf{Output:} Traceable knowledge, testable physical models, and recoverable research records.
\end{frame}

% Source page: Deepfield.md Page 56; exactly one frame.
\setsource{}
\begin{frame}[t,label=md-56]{II. Hierarchical Scientific Hypothesis Generation}
\small
\begin{itemize}
\setlength{\itemsep}{0.55em}
    \item \textbf{Ask across scales:} Nuclear structure and measurements through networks, astrophysics, and fundamental physics.
    \item \textbf{Transfer ideas:} Problem $\rightarrow$ challenge $\rightarrow$ gap $\rightarrow$ intuition $\rightarrow$ proposal.
    \item \textbf{Generate competitors:} Specify mechanisms, observable predictions, and falsification conditions.
    \item \textbf{Calculate and revise:} Use Code Agent scans, sensitivity, uncertainty, and feedback from both experiment types.
    \item \textbf{Evaluate:} Novelty, plausibility, information value, testability, feasibility, and complementarity.
\end{itemize}

\medskip
\textbf{Output:} Versioned hypothesis cards with evidence and quantitative tests.
\end{frame}

% Source page: Deepfield.md Page 57; exactly one frame.
\setsource{\href{https://doi.org/10.1038/s41586-026-10658-6}{ERA, Nature (2026)}}
\begin{frame}[t,label=md-57]{III. Code Reuse, Development, Analysis, and Simulation}
\small
\slideheading{Reuse first}
\smallskip

Search $\rightarrow$ reproduce $\rightarrow$ check validity $\rightarrow$ reuse
$\rightarrow$ adapt or extend $\rightarrow$ fill remaining gaps $\rightarrow$ validate.
\medskip
\begin{itemize}
\setlength{\itemsep}{0.5em}
    \item Prioritize existing reaction, network, stellar, and analysis software.
    \item Use execution feedback and code search for necessary extensions.
    \item Extract cross sections and rates with calibration, background, efficiency, and uncertainty checks.
    \item After measurement analysis, \textbf{proactively propose follow-up observations or simulations}.
    \item Run astrophysical and counterfactual calculations; archive code, environments, provenance, and results.
\end{itemize}
\end{frame}

% Source page: Deepfield.md Page 58; exactly one frame.
\setsource{\href{https://doi.org/10.1103/PhysRevAccelBeams.25.044601}{Miskovich et al., PRAB 25, 044601 (2022)}; \href{https://doi.org/10.1103/PhysRevResearch.7.013074}{Tsintari et al., PRR 7, 013074 (2025)}}
\begin{frame}[t,label=md-58]{IV. Experimental Design and Optimization}
\small
\begin{itemize}
\setlength{\itemsep}{0.65em}
    \item \textbf{Formulate:} Observables, competing hypotheses, uncertainties, constraints, and information value.
    \item \textbf{Design offline:} Optimize energies, beam, target, detectors, time, and resource allocation.
    \item \textbf{Simulate:} Background, efficiency, response, feasibility, and objective trade-offs.
    \item \textbf{Adapt online:} Bayesian updates and optimization using measurements and scientist feedback.
    \item \textbf{Prepare execution:} Validate interfaces in simulation or replay; coordinate follow-up observations.
\end{itemize}

\medskip
\textbf{Output:} Executable plans, configurations, calibration, and adjustment strategies; Direction V executes them.
\end{frame}

% Source page: Deepfield.md Page 59; exactly one frame.
\setsource{}
\begin{frame}[t,label=md-59]{V. Scientific Experimentation}
\small
\begin{itemize}
\setlength{\itemsep}{0.9em}
    \item \textbf{Beam experiments:} Real cross-section measurements, calibration, backgrounds, and repetitions with the experimental team.
    \item \textbf{Astrophysical experiments:} Targeted observations or independent archival analyses that test nuclear-constrained predictions.
    \item \textbf{Numerical experiments:} Reuse reaction and stellar codes for controlled variations, model comparisons, and uncertainty propagation.
\end{itemize}

\medskip
Beam measurement is the first experiment type; the post-measurement
observation or simulation is the second.
\medskip

{\footnotesize Preserve records, negative results, and unresolved explanations.
Distinguish simulation deductions from independent observational evidence.}
\end{frame}

